\documentclass[a4paper]{article}
\usepackage{amsfonts}
\usepackage{amsmath}
\usepackage{amssymb}
\usepackage{graphicx}
\usepackage{extarrows}

\newcommand{\limit}{\lim\limits}

\begin{document}

\noindent \large Auteur: Seda Jusjajeva\\
\noindent \large Studierichting: Informatica\\
\noindent \large Oefeningenreeks 3F nr. 1.9 \\

\medskip

\normalsize


$f(x) = \dfrac{5x^2 - x - 1}{x - 4}$ \\

VA: \\

$\limit_{x\to4}\dfrac{5x^2 - x - 1}{x - 4} = \dfrac{75}{0} =\infty\Rightarrow$ VA: $x = 4$ \\

$\limit_{x\to4+}\dfrac{5x^2 - x - 1}{x - 4} = \dfrac{75}{0^+} = +\infty$ \\

$\limit_{x\to4-}\dfrac{5x^2 - x - 1}{x - 4} = \dfrac{75}{0^-} = -\infty$ \\

HA: \\

$\limit_{x\to\infty}\dfrac{5x^2 - x - 1}{x - 4} = \limit_{x\to\infty}\dfrac{5x^2}{x} = \infty\Rightarrow$ geen HA \\

SA: \\

$\limit_{x\to\infty}\dfrac{f(x)}{x} = \limit_{x\to\infty}\dfrac{5x^2 - x - 1}{x(x - 4)} = \limit_{x\to\infty}\dfrac{5x^2}{x^2} = 5$ \\

$\limit_{x\to\infty}\left(f(x) - 5x\right) = \limit_{x\to\infty}\left(\dfrac{5x^2 - x - 1}{x - 4} - 5x\right)$ \\

$= \limit_{x\to\infty}\dfrac{5x^2 - x - 1 - 5x^2 + 20x}{x - 4}$ \\

$= \limit_{x\to\infty}\dfrac{19x}{x} = 19\Rightarrow$ SA: $y = 5x + 19$\\

$\dfrac{5x^2 - x - 1}{x - 4} - (5x + 19) = \dfrac{75}{x - 4}$ \\

$\begin{array}{r|ccc}
    x                 &   & 4     &   \\
    \hline  
    \dfrac{75}{x - 4} & - & \vert & + \\
\end{array}$ \\

$\Rightarrow$ Als $x\rightarrow+\infty\Rightarrow f(x) > A$ \\

$\Rightarrow$ Als $x\rightarrow-\infty\Rightarrow f(x) < A$ \\



\begin{thebibliography}{999}
%\bibitem{a} Referentie
\end{thebibliography}
\end{document}